\section{Divisor obstruction and Route-A assessment}

The clock theorem already stops the proposed closed Euler product.  The
scattering coefficient also has a global analytic obstruction to being one
entire Riemann target under source-allowed normalizations.

Define
\[
 \Lambda(u)=\pi^{-u/2}\Gamma(u/2)\zeta(u),\qquad
 \xi(u)=\tfrac12u(u-1)\Lambda(u).
\]
Then
\begin{equation}\label{eq:phi-xi}
 \PhiSc(s)=\frac{s}{s-1}\frac{\xi(2s-1)}{\xi(2s)}.
\end{equation}

\begin{theorem}[Zero-free normalization cannot remove the shifted divisor]
\label{thm:divisor}
Let $a,b\in\C$ with $a\ne0$, and let $h$ be entire and zero-free.  Then
$h(s)\PhiSc(as+b)$ is not entire.  In particular, it cannot equal the entire
function $\xi(s)$ as a global meromorphic identity.
\end{theorem}

\begin{proof}
Let $\rho$ be a nontrivial zero of $\zeta$, of multiplicity $m$.  At
$s=\rho/2$, the denominator $\Lambda(2s)$ vanishes.  The numerator is
$\Lambda(\rho-1)=\Lambda(2-\rho)$ by the functional equation.  Since
$\Re(2-\rho)>1$, this value is nonzero.  Thus $\PhiSc$ has a pole of
multiplicity $m$ at $\rho/2$.

At $s=(1+\rho)/2$, the numerator vanishes and the denominator
$\Lambda(1+\rho)$ is nonzero, so $\PhiSc$ has a shifted zero.  The affine map
$s\mapsto as+b$ is bijective, and an entire zero-free $h$ is finite and
nonzero at every transported pole.  Hence multiplication by $h$ cannot
remove the poles.
\end{proof}

A cusp-coordinate change $y'=ry$, followed by renormalization of the incoming
coefficient to one, sends $\PhiSc(s)$ to
$r^{2s-1}\PhiSc(s)$.  The extra exponential is entire and zero-free, so it
lies inside \cref{thm:divisor}.  A compensator with its own zeros or a
meromorphic zeta divisor could cancel poles, but would introduce arithmetic
not supplied by a zero-free cusp normalization.  The theorem is global and
does not prohibit a local identity on a pole-free domain.

The resulting Route-A decision is deliberately negative for the candidate
and positive for the obstruction package.

\begin{center}
\begin{tabular}{@{}lll@{}}
\toprule
Gate & Verdict & Decisive reason \\
\midrule
A1 primitive object & fail & open channel $\ne$ closed class \\
A2 same-clock determinant & fail & repetition forces $F=0$ \\
A3 global analytic target & fail & two shifted divisors \\
A4 operator lift & not testable & no same-clock operator \\
\bottomrule
\end{tabular}
\end{center}

The physical-line scattering matrix is nevertheless a natural unitary
operator-valued object.  That classical structure does not turn its resonant
divisor into a discrete self-adjoint Hilbert--P\'olya point spectrum.  Route B
is therefore not invoked.
